\cleardoublepage

\chapter{Anexos I}
\label{anexos-i}

Este anexo reúne material que no es de elaboración propia, sino tomado de otros autores, pero que respalda algunas decisiones del trabajo: la comparación de arquitecturas que justifica el uso de una TCN y las ventanas de \textit{declustering} aplicadas al catálogo.

\section{Comparativa empírica TCN vs redes recurrentes}
\label{anexo-tcn-lstm}

La Tabla~\ref{tab:tcn-vs-lstm} reproduce los resultados de \citet{bai2018tcn} sobre la comparación sistemática entre \emph{Temporal Convolutional Networks} (TCN), \emph{Long Short-Term Memory} (LSTM), \emph{Gated Recurrent Units} (GRU) y redes recurrentes simples (RNN) en una batería heterogénea de tareas de modelado de secuencias: tareas sintéticas (\textit{adding problem}, \textit{copy memory}), reconocimiento de imagen secuencializado (\textit{Sequential} y \textit{Permuted MNIST}), modelado polifónico de música (\textit{JSB Chorales}, \textit{Nottingham}) y modelado de lenguaje a nivel de palabra y carácter (\textit{Penn TreeBank}, \textit{Wikitext-103}, \textit{LAMBADA}, \textit{text8}).

\begin{table}[ht!]
\centering
\resizebox{\textwidth}{!}{%
\begin{tabular}{@{}llrrrr@{}}
\toprule
\textbf{Tarea} & \textbf{Métrica} & \textbf{LSTM} & \textbf{GRU} & \textbf{RNN} & \textbf{TCN} \\
\midrule
Sequential MNIST          & Accuracy (\%) $\uparrow$ & 87.2          & 96.2            & 21.5   & \textbf{99.0}   \\
Permuted MNIST            & Accuracy (\%) $\uparrow$ & 85.7          & 87.3            & 25.3   & \textbf{97.2}   \\
Adding problem ($T=600$)  & Loss $\downarrow$        & 0.164         & \textbf{5.3e-5} & 0.177  & 5.8e-5          \\
Copy memory ($T=1000$)    & Loss $\downarrow$        & 0.0204        & 0.0197          & 0.0202 & \textbf{3.5e-5} \\
Music JSB Chorales        & NLL $\downarrow$         & 8.45          & 8.43            & 8.91   & \textbf{8.10}   \\
Music Nottingham          & NLL $\downarrow$         & 3.29          & 3.46            & 4.05   & \textbf{3.07}   \\
Word-level Penn TreeBank  & Perplexity $\downarrow$  & \textbf{78.93} & 92.48          & 114.50 & 88.68           \\
Word-level Wikitext-103   & Perplexity $\downarrow$  & 48.4          & n/a             & n/a    & \textbf{45.19}  \\
LAMBADA                   & Perplexity $\downarrow$  & 4186          & n/a             & n/a    & \textbf{1279}   \\
Char-level Penn TreeBank  & bpc $\downarrow$         & 1.36          & 1.37            & 1.48   & \textbf{1.31}   \\
Char-level text8          & bpc $\downarrow$         & \textbf{1.41} & 1.42            & 1.69   & 1.45            \\
\bottomrule
\end{tabular}%
}
\caption[Comparación empírica TCN vs RNN/LSTM/GRU en \textit{benchmarks} de modelado de secuencias (\citealp{bai2018tcn})]{\textbf{Comparación empírica TCN vs RNN/LSTM/GRU en \textit{benchmarks} de modelado de secuencias.} Reproducido de la Tabla~1 de \citet{bai2018tcn}. Los valores en negrita indican el mejor resultado por tarea. La flecha $\uparrow$ indica que valores más altos son mejores; $\downarrow$, que valores más bajos son mejores. \texttt{n/a} cuando el modelo no fue evaluado por los autores. TCN obtiene el mejor resultado en 8 de las 11 tareas, frente a 2 de LSTM y 1 de GRU.}
\label{tab:tcn-vs-lstm}
\end{table}


Este estudio empírico es la principal evidencia citada en la literatura para justificar el uso de TCN como reemplazo de las arquitecturas recurrentes en problemas de series temporales, y motiva su elección en este TFM como columna vertebral del modelo descrito en el capítulo~\ref{desarrollo-del-proyecto-y-resultados}.

\section{Ventanas de \textit{declustering} de Gardner--Knopoff}
\label{anexo-gardner-knopoff}

El \textit{declustering} del catálogo (Sección~\ref{declustering}) emplea las ventanas
espacio-temporales de \citet{gardner1974poissonian}, en la formulación continua de
\citet{vanstiphout2012declustering}. Para cada evento candidato a \textit{mainshock} de magnitud
\(M\) se define un radio espacial \(L(M)\)~\eqref{eq:gk-L} y una duración temporal \(T(M)\)~\eqref{eq:gk-T},
y se descartan como dependientes todos los eventos de menor magnitud que caen dentro de esa ventana.
La Tabla~\ref{tab:gk-ventanas} muestra los valores de ambas ventanas para varias magnitudes; el
salto en \(T(M)\) a partir de \(M = 6.5\) corresponde al cambio de rama de la fórmula original.

\begin{table}[ht!]
  \centering
  \small
  \begin{tabular}{@{}ccc@{}}
    \toprule
    \textbf{Magnitud \(M\)} & \textbf{Radio \(L(M)\) (km)} & \textbf{Duración \(T(M)\) (días)} \\
    \midrule
    4.5                     & 34.7                         & 77                                \\
    5.0                     & 40.0                         & 144                               \\
    5.5                     & 46.1                         & 268                               \\
    6.0                     & 53.2                         & 499                               \\
    6.5                     & 61.3                         & 885                               \\
    7.0                     & 70.7                         & 918                               \\
    7.5                     & 81.6                         & 953                               \\
    \bottomrule
  \end{tabular}
  \caption[Ventanas de \textit{declustering} de Gardner--Knopoff]{\textbf{Ventanas de Gardner--Knopoff} en su forma continua \citep{vanstiphout2012declustering}: radio espacial \(L(M) = 10^{0.1238\,M + 0.983}\) km y duración temporal \(T(M)\), que cambia de rama en \(M = 6.5\).}
  \label{tab:gk-ventanas}
\end{table}
